\subsection{Reactive Type System of \GRFRP}
\label{sec:semantics:grfrp:react_typing}

As stated in \Cref{sec:semantics:grfrp:syntax}, components invoked within services must be
\emph{reactive}. This static property, introduced in \Cref{sec:semantics:grsync:react_typing},
tracks dependencies on \emph{continuous flows}. A continuous flow is one whose evolution depends
on the execution rate itself, rather than on discrete input changes. Such flows are
problematic for an event-driven model, which assumes computations are triggered only by new
input values. We identify three primary sources of this continuity.

The first source is the \timeop operation, which returns the current execution time. Its value
changes continuously, an observation will always give a value different from the last observation.

The second arises from an untriggered feedback loop using a \kwf{last} memory. For instance,
consider the simple counter component from \Cref{sec:semantics:grsync:react_typing}:
\begin{equation*}
    \component{\idf{count}}{}{\idf{o}}{\idf{o} = \last{\idf{o}} + 1}{\idf{o} = 0}.
\end{equation*}
From an event-driven perspective, this component should remain inactive, as it has no inputs
to trigger a computation. However, its definition implies a self-perpetuating loop where the
value of \idf{o} depends on its own previous value, \last{\idf{o}}. The flow is thus
\emph{continuous}, as it should be continuously re-evaluated in order to maintain its value
up-to-date. This violates the principles of reactive, event-driven execution.

The third source is the untriggered emission of an event with the \kwf{emit} operator.
Consider the \idf{emit\_now} component, also from \Cref{sec:semantics:grsync:react_typing}:
\begin{equation*}
    \component{\idf{emit\_now}}{\idf{s}}{\idf{e}}{\idf{e} = \emit{\idf{s}}}{}.
\end{equation*}
Conceptually, the output event \idf{e} is \emph{continuously} present. However, an
event-driven execution would trigger the component only when the input signal \idf{s}
changes, emitting the event only at that moment. This mismatch between the intended
continuous presence and the discrete nature of event-driven propagation creates ambiguity.

To resolve these conflicts, the use of \timeop or untriggered \kwf{last} memories and event
emissions must be restricted. We formalize this restriction through a reactive type system that
distinguishes between two kinds of flows:
%
\begin{itemize}
    \item \emph{Continuous flows}, which originate from sources inherently tied to the
          execution step itself, such as \timeop or untriggered memories and event emissions.
    \item \emph{Discrete flows}, which encompass all other flows. These are driven by input
          changes and are stable between events: signals hold their last value, and events
          are absent.
\end{itemize}

This distinction is enforced statically by the reactive type system.
\Cref{sec:semantics:grsync:react_typing} introduced the \GRSync portion of this system, which
defines typing rules to qualify a component as reactive. This section presents the \GRFRP
portion. Enforcing this type system for services enables the oversampling theorem
(\Cref{thm:semantics:grfrp:theorems:oversampling}). This theorem guarantees that propagating
input changes is sufficient to keep outputs up-to-date, and that re-computing them without
new inputs does not alter their values.

\begin{table}[!ht]
    \centering
    \begin{tabularx}{\textwidth}{|l|X|}\hline
        \rType ::= \cType $\vert$ \dType        & Reactive type denoting if the
        associated element is continuous (depending on \timeop or untriggered event emissions and
        \kwf{last} memories) or discrete.                                                     \\\hline
        $\Theta[x] = \rType$                    & A reactive typing context $\Theta$
        mapping flow identifiers to their reactive type.                                      \\\hline
        \servTypedReact{F}                      & The service $F$ is reactive.                \\\hline
        \stmtTypedReact{\Theta}{stmt^+}         & The statements $stmt^+$ are reactive
        in the reactive typing context $\Theta$.                                              \\\hline
        \opTypedReact{\Theta}{\sigma}{\rType^+} & The flow operation $\sigma$ is reactive and
        has type $\rType^+$ in the reactive typing context $\Theta$.                          \\\hline
    \end{tabularx}
    \caption[Notations for the reactive type system of \GRFRP.]%
    {Notations for the logic of the reactive type system of \GRFRP.}
    \label{tab:semantics:grfrp:react_typing:notations}
\end{table}

The notations for the reactive type system, which recall and extend those for \GRSync from
\Cref{sec:semantics:grsync:react_typing}, are summarized in
\Cref{tab:semantics:grfrp:react_typing:notations}.
%
The reactive type \rType detects dependencies on continuous sources. A typing context $\Theta$
maps each flow identifier $x$ to its type. $\Theta[x] = \cType$ denotes that $x$ is a
\emph{continuous} flow, while $\Theta[x] = \dType$ denotes that it is a \emph{discrete} flow.
%
Typing judgments capture reactivity at different levels of the syntax. \servTypedReact{F}
asserts that a service $F$ is reactive. \stmtTypedReact{\Theta}{stmt^+} asserts that a
sequence of statements is reactive in a context $\Theta$. Finally,
\opTypedReact{\Theta}{\sigma}{\rType^+} asserts that a flow operation $\sigma$ has type
$\rType^+$ in context $\Theta$.

The following paragraphs present the rules that define the \GRFRP portion of this reactive
type system.

\paragraph{Services}

A service is reactive if its implementation adheres to the constraints of the type system,
ensuring its soundness under an event-driven execution model.
%
According to \nameref{rule:semantics:grfrp:react_typing:serv}, a service
\service{F}{stmt^+} is reactive if there exists a typing context $\Theta$ under which its
block of statements $stmt^+$ is reactive.

\begin{center}
    \begin{minipage}{0.95\textwidth}
        \begin{mathpar}
            \inferREACT
            {
                G[F] = \service{F}{stmt^+} \\
                \stmtTypedReact{\Theta}{stmt^+}
            }
            { \servTypedReact{F} }
            {Serv}
            \label{rule:semantics:grfrp:react_typing:serv}
        \end{mathpar}
    \end{minipage}
\end{center}

\paragraph{Statements}
%
Statements propagate reactive types from flow operations to identifiers.
%
As formalized by \nameref{rule:semantics:grfrp:react_typing:stmts}, a block of statements
($stmt^+$) is reactive if each individual statement ($stmt$) is reactive.

\begin{center}
    \begin{minipage}{0.95\textwidth}
        \begin{mathpar}
            \inferREACTFlowStmt
            { \left(\stmtTypedReact{\Theta}{stmt}\right)^+ }
            { \stmtTypedReact{\Theta}{stmt^+} }
            {Blck}
            \label{rule:semantics:grfrp:react_typing:stmts}
        \end{mathpar}
    \end{minipage}
\end{center}

As stated by \nameref{rule:semantics:grfrp:react_typing:stmt:decl}, a declaration
(\decl{x^+}{\sigma}) is reactive if the flow operation $\sigma$ is reactive and its resulting
types match the types assigned to $x^+$ in the context $\Theta$.

\begin{center}
    \begin{minipage}{0.95\textwidth}
        \begin{mathpar}
            \inferREACTFlowStmt
            { \opTypedReact{\Theta}{\sigma}{(\Theta[x])^+} }
            { \stmtTypedReact{\Theta}{\decl{x^+}{\sigma}} }
            {Decl}
            \label{rule:semantics:grfrp:react_typing:stmt:decl}
        \end{mathpar}
    \end{minipage}
\end{center}

The rules \nameref{rule:semantics:grfrp:react_typing:stmt:import} and
\nameref{rule:semantics:grfrp:react_typing:stmt:export} for \kwf{import} and \kwf{export} require
the associated flows to be discrete (of type \dType).
This enforces a clear boundary for the service, ensuring its external interface only involves
discrete flows and preventing continuous behaviors from leaking in or out.

\begin{center}
    \begin{minipage}{0.95\textwidth}
        \begin{mathpar}
            \inferREACTFlowStmt
            { \Theta[x] = \dType }
            { \stmtTypedReact{\Theta}{\import{x}} }
            {Import}
            \label{rule:semantics:grfrp:react_typing:stmt:import}

            \inferREACTFlowStmt
            { \Theta[x] = \dType }
            { \stmtTypedReact{\Theta}{\export{x}} }
            {Export}
            \label{rule:semantics:grfrp:react_typing:stmt:export}
        \end{mathpar}
    \end{minipage}
\end{center}

\paragraph{Flow operations}

The typing rules for flow operations determine if a resulting flow is continuous or discrete,
thereby propagating dependencies on continuous sources.
%
An identifier ($x$) has the type assigned to it in the context $\Theta$, as formalized by
\nameref{rule:semantics:grfrp:react_typing:op:ident}.

\begin{center}
    \begin{minipage}{0.95\textwidth}
        \begin{mathpar}
            \inferREACTFlow
            { }
            { \opTypedReact{\Theta}{x}{\Theta[x]} }
            {ID}
            \label{rule:semantics:grfrp:react_typing:op:ident}
        \end{mathpar}
    \end{minipage}
\end{center}

The \timeop operation is the primary source of continuity and is therefore always of type
\cType, as per \nameref{rule:semantics:grfrp:react_typing:op:time}.

\begin{center}
    \begin{minipage}{0.95\textwidth}
        \begin{mathpar}
            \inferREACTFlow
            { }
            { \opTypedReact{\Theta}{\timeop}{\cType} }
            {Time}
            \label{rule:semantics:grfrp:react_typing:op:time}
        \end{mathpar}
    \end{minipage}
\end{center}

A component application $f(\sigma^+)$ is reactive if the component $f$ is itself reactive,
given the types of its inputs. As specified by
\nameref{rule:semantics:grfrp:react_typing:op:app}, this check relies on the component typing
judgment from \GRSync.
The reactivity of a component is crucial to ensure a sound usage in an asynchronous service.
\Cref{thm:semantics:grsync:theorems:oversampling} states that a reactive component can be executed
in an event-driven manner, only propagating changes. This execution is proved to be resilient to
oversampling: the absence of changes from the discrete input flows of the component leads to an
absence of changes from its internal memory and from its discrete output flows.

\begin{center}
    \begin{minipage}{0.95\textwidth}
        \begin{mathpar}
            \inferREACTFlow
            {
                \left(\opTypedReact{\Theta}{\sigma}{\rType_i}\right)^+ \\
                \compTypedReact{f}{\rType_i^+}{\rType_o^+}
            }
            { \opTypedReact{\Theta}{f(\sigma^+)}{\rType_o^+} }
            {App}
            \label{rule:semantics:grfrp:react_typing:op:app}
        \end{mathpar}
    \end{minipage}
\end{center}

All events must occur at precise, discrete moments. The \onchange{\sigma} and
\mergeop{\sigma_1}{\sigma_2} operations produce events and thus require their inputs to be
discrete, ensuring their outputs are also discrete.

\begin{center}
    \begin{minipage}{0.95\textwidth}
        \begin{mathpar}
            \inferREACTFlow
            { \opTypedReact{\Theta}{\sigma}{\dType} }
            { \opTypedReact{\Theta}{\onchange{\sigma}}{\dType} }
            {OChg}
            \label{rule:semantics:grfrp:react_typing:op:onchange}

            \inferREACTFlow
            { \forall i \in \{1,2\}, \; \opTypedReact{\Theta}{\sigma_i}{\dType} }
            { \opTypedReact{\Theta}{\mergeop{\sigma_1}{\sigma_2}}{\dType} }
            {Mrg}
            \label{rule:semantics:grfrp:react_typing:op:merge}
        \end{mathpar}
    \end{minipage}
\end{center}

Similarly, the \persist{\sigma} and \sample{\sigma}{T} operations are triggered by an event
$\sigma$, as specified by
\nameref{rule:semantics:grfrp:react_typing:op:persist} and
\nameref{rule:semantics:grfrp:react_typing:op:sample}. To ensure the output is produced at a precise
time, the input event must be discrete, resulting in a discrete output.

\begin{center}
    \begin{minipage}{0.95\textwidth}
        \begin{mathpar}
            \inferREACTFlow
            { \opTypedReact{\Theta}{\sigma}{\dType} }
            { \opTypedReact{\Theta}{\persist{\sigma}}{\dType} }
            {Prst}
            \label{rule:semantics:grfrp:react_typing:op:persist}

            \inferREACTFlow
            { \opTypedReact{\Theta}{\sigma}{\dType} }
            { \opTypedReact{\Theta}{\sample{\sigma}{T}}{\dType} }
            {Samp}
            \label{rule:semantics:grfrp:react_typing:op:sample}
        \end{mathpar}
    \end{minipage}
\end{center}

The \scan{\sigma}{T} operation produces periodic samples of the signal $\sigma$, as formalized by
rule \nameref{rule:semantics:grfrp:react_typing:op:scan}. Because the sampling is driven by a
fixed-period timer, the output is always a discrete signal, regardless of the input type. This
operation is one of the few ways to discretize a continuous flow.

\begin{center}
    \begin{minipage}{0.95\textwidth}
        \begin{mathpar}
            \inferREACTFlow
            { \opTypedReact{\Theta}{\sigma}{\rType} }
            { \opTypedReact{\Theta}{\scan{\sigma}{T}}{\dType} }
            {Scan}
            \label{rule:semantics:grfrp:react_typing:op:scan}
        \end{mathpar}
    \end{minipage}
\end{center}

The \timeout{\sigma}{T} operation creates a timer that resets on an input event. The input
event must be discrete to ensure that the reset occurs at a precise moment, as stated by
\nameref{rule:semantics:grfrp:react_typing:op:timeout}. If the input timing were imprecise, the
reset will also be imprecise, and all subsequent timeout events would also have imprecise timing.

\begin{center}
    \begin{minipage}{0.95\textwidth}
        \begin{mathpar}
            \inferREACTFlow
            { \opTypedReact{\Theta}{\sigma}{\dType} }
            { \opTypedReact{\Theta}{\timeout{\sigma}{T}}{\dType} }
            {TOut}
            \label{rule:semantics:grfrp:react_typing:op:timeout}
        \end{mathpar}
    \end{minipage}
\end{center}

The \throttle{\sigma}{\Delta} operation filters a signal. If the input signal $\sigma$ were
continuous, the exact moment its value crosses the $\Delta$ threshold would be ill-defined in
an event-driven model. To avoid this ambiguity, the input must be discrete, as stated by the rule
\nameref{rule:semantics:grfrp:react_typing:op:throttle}, which in turn
makes the output discrete.

\begin{center}
    \begin{minipage}{0.95\textwidth}
        \begin{mathpar}
            \inferREACTFlow
            { \opTypedReact{\Theta}{\sigma}{\dType} }
            { \opTypedReact{\Theta}{\throttle{\sigma}{\Delta}}{\dType} }
            {Thrt}
            \label{rule:semantics:grfrp:react_typing:op:throttle}
        \end{mathpar}
    \end{minipage}
\end{center}

\paragraph{Summary}

This section has detailed the reactive type system for \GRFRP, which complements the system for
\GRSync components. By classifying flows as either \emph{continuous} or \emph{discrete}, the type
system imposes constraints on service definitions. It mandates that service interfaces (\ie imports
and exports) and most event-based operations handle only discrete flows. This disciplined approach
prevents the propagation of continuous behaviors that are incompatible with an event-driven
execution model. Consequently, any well-typed service is guaranteed to be reactive, establishing a
formal foundation for the correctness of its asynchronous execution
(\cf \Cref{thm:semantics:grfrp:theorems:oversampling}).
